\documentclass[11pt,letterpaper]{article}

\usepackage[margin=1in]{geometry}
\usepackage{amsmath,amssymb}
\usepackage{booktabs}
\usepackage{enumitem}
\usepackage{xcolor}
\usepackage{hyperref}
\usepackage{titlesec}
\usepackage{microtype}
\usepackage{parskip}
\usepackage{longtable}
\usepackage{array}
\usepackage{tabularx}

\hypersetup{
  colorlinks=true,
  linkcolor=blue!50!black,
  urlcolor=blue!50!black,
  pdftitle={Proof Brokerage: Implementation Roadmap}
}

\definecolor{phasecolor}{rgb}{0.1, 0.3, 0.6}
\definecolor{riskhigh}{rgb}{0.7, 0.1, 0.1}
\definecolor{riskmed}{rgb}{0.7, 0.5, 0.1}
\definecolor{risklow}{rgb}{0.1, 0.5, 0.1}

\titleformat{\section}[hang]{\normalfont\Large\bfseries\color{phasecolor}}{\thesection}{0.6em}{}
\titleformat{\subsection}[hang]{\normalfont\large\bfseries}{\thesubsection}{0.6em}{}

\newcommand{\highrisk}{\textcolor{riskhigh}{\textbf{HIGH}}}
\newcommand{\medrisk}{\textcolor{riskmed}{\textbf{MED}}}
\newcommand{\lowrisk}{\textcolor{risklow}{\textbf{LOW}}}

\title{
  \vspace{-1cm}
  \textbf{Proof Brokerage: Implementation Roadmap}\\[0.3em]
  \large{From spec v1.0 to working v1 system}
}
\author{}
\date{Working draft --- companion to spec v1.0}

\begin{document}
\maketitle

\sloppy
\emergencystretch=3em

\begin{abstract}
\noindent
This document is the implementation companion to the Proof Brokerage
Architecture Specification v1.0. It sequences the engineering work into
phases, identifies dependencies between components, surfaces the highest
risks, and provides milestone-level deliverables. It is explicitly a
\emph{plan}, not a contract: the sequencing reflects what the
specification authors believe is the right order to build, and is
expected to be revised as implementation surfaces information not
available at design time.

The roadmap targets a 12--18 month timeline to a working v1 system.
Critical-path components (the IR, the Lean plugin, CVC5 integration)
are front-loaded; lower-priority work (Rocq probe, LLM integration, the
dashboard) is deferred until the critical path validates. The roadmap
does not address staffing or organizational structure, which are
project-specific.
\end{abstract}

\tableofcontents
\newpage

\section{Roadmap philosophy}

The architecture has many components. Implementing them in the wrong
order risks substantial wasted work. Two principles drive the
sequencing:

\paragraph{Validate the IR before building on it.} The IR is the
load-bearing artifact of the architecture; if it has design flaws, every
downstream component built against it has to be revised. The earliest
phases focus on validating the IR with one home system and one backend,
end-to-end, before adding more.

\paragraph{Defer LLM and dashboard work until the symbolic core works.}
Both LLM integration and the dashboard are genuinely valuable, but they
can be added once the symbolic-only core is operational. Building them
in parallel risks the LLM components masking weaknesses in the symbolic
path that we'd rather know about early.

The roadmap is structured as five phases. Phases are sequential at the
component level; within a phase, components can be developed in
parallel where dependencies permit.

\section{Phase summary}

\begin{tabularx}{\textwidth}{@{}llXr@{}}
\toprule
\textbf{Phase} & \textbf{Codename} & \textbf{Outcome} & \textbf{Est.\ duration} \\
\midrule
0 & Foundations  & Repository, schemas, tooling                             & 1 month \\
1 & Skeleton     & End-to-end Tier 1 dispatch, Lean + CVC5, simplest goals  & 3 months \\
2 & Core         & Tier 3 with Alethe; full IR rewriter; metadata richness  & 4 months \\
3 & Breadth      & Vampire (ATP) + LLM-as-backend; concurrent dispatch       & 3 months \\
4 & Probe        & Rocq architectural probe (Tier 0+1)                      & 2 months \\
5 & Polish       & Dashboard, build path, caching, Tier 2 framing           & 2 months \\
\bottomrule
\end{tabularx}

\bigskip

Total estimated duration: 15 months. Phase 5 may overlap with Phase 4
where independent.

\section{Phase 0: Foundations}
\label{sec:phase0}

\paragraph{Goal.} Establish the repository structure, the canonical
schemas, the development tooling, and the cross-component testing
harness. Nothing in this phase produces a working dispatch; everything
in this phase is prerequisite to phases that do.

\subsection{Deliverables}

\begin{enumerate}
  \item \textbf{Schema repository.} JSON-Schema definitions for the IR
        document, the certificate envelope, the refinement record, the
        adapter capability manifest, and the rewrite trace. Versioned
        from v1.0; CI validates that schema changes don't break
        round-trip with sample documents.
  \item \textbf{Registry document.} The patterns registry (logic
        features, fragments, theory tags, axiom shape descriptors,
        concept tags, construction kinds) as a versioned data file.
        Initial population from spec Appendix A. Maintenance process
        documented.
  \item \textbf{Reference IR documents.} Hand-written IR documents for
        the three worked examples in spec Section 12. These are the
        canonical test inputs that all components must handle.
  \item \textbf{Cross-language SDK skeleton.} Language bindings for
        Lean (in Lean), for Rocq (in OCaml), and a reference
        implementation library (in Rust or OCaml) shared by
        adapters. Decided in this phase, not before.
  \item \textbf{Continuous integration.} Build and test pipelines for
        all components, with fast cycle times.
\end{enumerate}

\subsection{Risks}

\begin{itemize}
  \item \medrisk\ Choice of implementation language for the shared
        library affects integration complexity for the home-system
        plugins. Recommendation: Rust for the shared library, with
        explicit FFI for Lean (via Lean's C interface) and for OCaml
        (via standard FFI). The cost of getting this wrong is rework
        in Phase 1.
  \item \lowrisk\ Schema drift between hand-written reference IRs
        and serializer outputs once the latter exists. Mitigation:
        treat reference IRs as canonical; serializers are validated
        against them.
\end{itemize}

\subsection{Exit criteria}

Schemas validate on all reference documents. Registry is populated and
documented. SDK skeleton compiles in all target languages. CI is green.

\section{Phase 1: Skeleton}
\label{sec:phase1}

\paragraph{Goal.} Working end-to-end dispatch from Lean 4 to CVC5 for
the simplest possible goal class: pure linear arithmetic, no
typeclasses, no quotients, no higher-order. Tier 1 (Farkas) only.
This phase is about validating that the architecture's pipeline works
at all; capability is intentionally minimal.

\subsection{Deliverables}

\begin{enumerate}
  \item \textbf{Lean plugin (minimal).} A \texttt{dispatch} tactic that
        serializes a proof state into the IR, restricted to goals of the
        form $\forall \vec{x} : \mathbb{Z}, P_1 \wedge \cdots \wedge P_n
        \to Q$ where each $P_i, Q$ is a linear arithmetic predicate.
        Other goals are rejected with a clear ``not yet supported''
        message.
  \item \textbf{IR rewriter (minimal).} Propositional simplification
        only. No quotient elimination, no unfolding.
  \item \textbf{Dispatcher (minimal).} Single backend, no capability
        matching, no concurrency.
  \item \textbf{CVC5 adapter (minimal).} Translates LIA shell into
        SMT-LIB. Receives unsat result. \emph{Does not yet consume
        Alethe traces}; produces a Tier 1 (Farkas) certificate by
        running CVC5 with the appropriate options.
  \item \textbf{Tier 1 verifier (Farkas).} Lean tactic that consumes a
        Farkas certificate and produces a kernel-checked proof term.
  \item \textbf{Lifting (trivial).} For Phase 1 goals, refinement
        records contain only \texttt{type\_specialization} from $\alpha$
        to $\mathbb{Z}$ (or none, for goals already over $\mathbb{Z}$),
        and lifting is mechanical.
\end{enumerate}

\subsection{Risks}

\begin{itemize}
  \item \highrisk\ The translation hints schema may be insufficient even
        for this restricted case. \emph{This is the single highest risk
        in the roadmap.} The worked examples in spec Section 12 are
        believed sufficient, but real Lean code will surface
        configurations not anticipated. Mitigation: budget 50\% of Phase
        1 time for IR/spec revision based on what the implementation
        surfaces. Treat the spec as a draft until Phase 1 completes.
  \item \medrisk\ Farkas certificate extraction from CVC5 may have
        format quirks. CVC5's certificate output has evolved; the
        adapter must handle the current version with explicit version
        guards.
  \item \medrisk\ Lean's metaprogramming API for accessing proof state
        is rich but intricate. The serializer is the most complex piece
        of the Lean plugin.
\end{itemize}

\subsection{Exit criteria}

The three worked examples from spec Section 12, restricted to the
LIA case (Example 1 with $\alpha = \mathbb{Z}$ explicitly), dispatch
end-to-end and produce kernel-checked proof terms. The Lean plugin
rejects out-of-scope goals with clear messages. Performance is acceptable
for interactive use (sub-10-second total dispatch latency on the
example goals).

\subsection{Spec revision checkpoint}

At Phase 1 exit, conduct an explicit specification review based on
implementation findings. Document any deviations from spec v1.0 and
update the spec to v1.1 if warranted. Phase 2 builds on the revised
spec, not on v1.0.

\section{Phase 2: Core}
\label{sec:phase2}

\paragraph{Goal.} Bring the IR and verifier to full v1 capability for
Lean and CVC5: typeclass-quantified goals, quotient types, the IR
rewriter with quotient elimination, Tier 3 with Alethe traces, full
lifting machinery.

\subsection{Deliverables}

\begin{enumerate}
  \item \textbf{Type metadata: full schema implementation.} The
        serializer produces type metadata for primitive types, type
        variables (with instances), and type constructor applications
        (quotient, inductive, sigma, refinement). Other construction
        kinds remain placeholder.
  \item \textbf{Definitional metadata: full schema implementation.} All
        kinds defined in the spec are implemented in the serializer.
  \item \textbf{Theory tag soundness mechanism.} The serializer
        validates theory tags against verified embedding lemmas. The
        registry is consulted at serialization time; tags without
        backing lemmas in the home system are dropped.
  \item \textbf{IR rewriter passes.} \texttt{quotient\_elimination},
        \texttt{definition\_unfolding} (selective by concept tag),
        \texttt{equality\_saturation}, \texttt{quantifier\_prenexing},
        \texttt{skolemization}. Each with rewrite trace recording for
        inversion.
  \item \textbf{CVC5 adapter: full.} Produces Tier 3 certificates with
        Alethe traces. Refinement records include all specialization
        kinds.
  \item \textbf{Alethe-to-Lean replayer.} Symbolic Tier 3 verifier
        that consumes Alethe traces and produces Lean proof terms.
        This is a substantial sub-project.
  \item \textbf{Lifting: full machinery.} Inverts refinement records
        and rewrite traces in the right order, handles all
        specialization kinds, handles all rewrite passes.
  \item \textbf{Capability matching.} Dispatcher consults adapter
        manifests against goal classification.
\end{enumerate}

\subsection{Risks}

\begin{itemize}
  \item \highrisk\ The Alethe-to-Lean replayer is itself a research-grade
        engineering project. Existing work (e.g., on Alethe-to-Isabelle)
        provides a starting point but cannot be transferred wholesale.
        Mitigation: scope the initial replayer to support only the
        Alethe rules CVC5 actually emits for the v1 backend
        configuration; expand coverage incrementally.
  \item \highrisk\ Rewrite trace inversion for \texttt{quotient\_elimination}
        is non-trivial. Getting it wrong produces silent unsoundness or
        loud lifting failures. Mitigation: extensive property-based
        testing on the inversion (round-trip: original goal $\to$
        rewritten $\to$ certificate $\to$ lifted = original goal).
  \item \medrisk\ Theory tag soundness validation requires a registry of
        verified embedding lemmas in mathlib. Some tags will need new
        lemmas to be proved. Estimate: 5--10 lemmas in mathlib, each a
        small but non-trivial proof.
  \item \medrisk\ Definitional metadata for typeclass methods has
        subtleties around instance resolution that are easy to get
        wrong. Phase 1's findings should inform this.
\end{itemize}

\subsection{Exit criteria}

All three worked examples from spec Section 12 dispatch end-to-end with
full capability. Performance is acceptable for interactive use on
representative real-world Lean goals (sub-30-second total dispatch
latency for typical mathlib-style goals). The Alethe replayer's coverage
is documented; goals outside coverage fall through to LLM-assisted
reconstruction (which is not yet implemented; in Phase 2 they fail).

\section{Phase 3: Breadth}
\label{sec:phase3}

\paragraph{Goal.} Add Vampire (ATP) and a generic LLM-as-backend
adapter. Implement concurrent dispatch. The architecture's heterogeneity
claims are validated only when multiple backends with different
characteristics are dispatching the same goal.

\subsection{Deliverables}

\begin{enumerate}
  \item \textbf{Vampire adapter.} TPTP-FOF and TPTP-THF input
        production. Refinement records with \texttt{axiomatization}
        specializations. Capability manifest declares first-order and
        higher-order coverage.
  \item \textbf{TPTP-to-Lean replayer.} Symbolic Tier 3 verifier for
        TPTP-FOF traces. May initially have lower coverage than Alethe;
        documented gaps fall through to LLM strategy.
  \item \textbf{LLM-as-backend adapter.} Renders IR as Lean surface
        syntax, prompts a configured LLM endpoint, packages the
        returned tactic script as a Tier 3 certificate.
  \item \textbf{LLM-assisted Tier 3 reconstruction strategy.} For
        certificates whose traces cannot be replayed symbolically, a
        configured LLM is asked to translate. Kernel-checked.
  \item \textbf{Concurrent dispatch.} Multiple backends in parallel,
        first-valid-wins with grace window for higher tiers, cancellation
        on success.
  \item \textbf{Adapter capability matching: refined.} Filtering and
        prioritization based on goal classification, with fallback
        chains.
\end{enumerate}

\subsection{Risks}

\begin{itemize}
  \item \medrisk\ Vampire-HOL's higher-order capabilities are powerful
        but uneven. Some higher-order goals will succeed in Vampire
        but fail to replay; this is acceptable for v1 but should be
        tracked.
  \item \medrisk\ LLM endpoint latency, cost, and availability are
        operational concerns that don't show up in Phase 2. The first
        real LLM use will surface practical issues (rate limiting,
        timeout calibration, prompt engineering).
  \item \medrisk\ The grace-window policy for higher-tier certificates
        in concurrent dispatch may need tuning. Default is 2 seconds
        but real workloads may want different defaults.
  \item \lowrisk\ Generic LLM prompts may underperform compared to
        specialized prover models (DeepSeek-Prover etc.). Mitigation:
        the LLM adapter accepts model identity as configuration; users
        can configure specialized models if available.
\end{itemize}

\subsection{Exit criteria}

A goal involving function composition (worked example 2) succeeds via
LLM-as-backend with sub-5-second latency and via Vampire-HOL within
30 seconds. Concurrent dispatch demonstrably picks the faster
successful backend. The dashboard (Phase 5) is not yet implemented;
diagnostics are command-line only.

\section{Phase 4: Probe}
\label{sec:phase4}

\paragraph{Goal.} Implement the Rocq architectural probe. The probe
validates that the IR is not Lean-shaped and that the protocol works
across home systems. Production-quality Rocq integration is a v2
goal; the probe is intentionally minimal.

\subsection{Deliverables}

\begin{enumerate}
  \item \textbf{Rocq plugin (probe scope).} A \texttt{Dispatch} tactic
        that serializes Rocq proof state into the IR. Coverage: the
        same goal class as Lean Phase 1 (LIA only, Tier 1) plus
        Phase 2 capability for the worked-example category.
  \item \textbf{Tier 0 verifier (Rocq).} Trust-policy lookup, oracle
        certificate acceptance with explicit trust expansion.
  \item \textbf{Tier 1 verifier (Rocq).} Farkas certificate consumption
        producing kernel-checked Rocq proof terms.
  \item \textbf{IR shape validation.} Run the Phase 2 reference IRs
        produced by the Lean serializer through the Rocq deserializer.
        Run the Phase 2 reference IRs produced by the Rocq serializer
        through the Lean adapter pipeline. Both directions must
        succeed (modulo Rocq probe scope restrictions). This is the
        validation that the IR is not Lean-shaped.
\end{enumerate}

\subsection{Risks}

\begin{itemize}
  \item \highrisk\ Rocq's typeclass system (canonical structures,
        type classes, Hint Cut, etc.) is genuinely different from
        Lean 4's. The IR may need extension. Mitigation: this is
        \emph{the} validation step for Position 3 of the architecture
        debate. If it fails badly, the architecture itself is wrong
        and Phases 2--3 may need partial revision.
  \item \medrisk\ Rocq's plugin development model (OCaml, MetaCoq,
        Ltac2) is more conservative than Lean's. Engineering effort
        per feature is meaningfully higher.
  \item \medrisk\ Cross-system reference IR validation will surface
        discrepancies that don't show up when each home system uses
        its own serializer. Some discrepancies will be benign
        (different but equivalent representations); others will reveal
        real schema gaps.
\end{itemize}

\subsection{Exit criteria}

Rocq probe dispatches simple LIA and worked-example-2-style goals
end-to-end. Cross-system IR round-trips for at least the three worked
examples succeed (with documented exceptions for cases the probe
explicitly does not support). Schema gaps discovered during this phase
are documented for v2 resolution.

\subsection{Spec revision checkpoint}

If Rocq probe finds material schema gaps, conduct another spec review.
This is the architectural-validation milestone. The spec moves to v1.2
or v2.0 depending on the severity of revisions.

\section{Phase 5: Polish}
\label{sec:phase5}

\paragraph{Goal.} Complete the user-facing experience: the dashboard,
the build path, the certificate cache, and Tier 2 with its UX
commitments. These are the components that turn a working pipeline
into a usable system.

\subsection{Deliverables}

\begin{enumerate}
  \item \textbf{Tier 2 implementation.} Lemma list extraction from
        backends that don't produce traces; reconstruction via Lean's
        \texttt{aesop} and \texttt{simp\_arith} with hints; structured
        failure reporting.
  \item \textbf{Tier 2 UX commitments.} Distinct labeling
        (\texttt{reconstructed (probabilistic)}); per-backend, per-goal-shape
        success rate tracking.
  \item \textbf{The dashboard.} Session and project-level dispatch
        history; per-backend, per-tier, per-goal-shape statistics;
        failure analysis with suggested remedies.
  \item \textbf{Build path integration.} Dispatch calls re-run at build
        time with persistent caching. CI integration documented.
  \item \textbf{Certificate cache: full.} Project-local store with all
        invalidation rules; optional shared-cache configuration.
  \item \textbf{Documentation and tutorials.} User-facing documentation,
        worked examples, troubleshooting guide.
\end{enumerate}

\subsection{Risks}

\begin{itemize}
  \item \medrisk\ The Tier 2 demotion clause may trigger. If Lean's
        \texttt{aesop}-with-hints reconstruction success rate is
        substantially worse than \texttt{sledgehammer}'s baseline, Tier 2
        is demoted to v1.5 and Phase 5 ships without it.
  \item \medrisk\ Dashboard design decisions are user-facing and may
        require iteration. Initial dashboard should be functional but
        may be redesigned based on early-user feedback.
  \item \lowrisk\ Build-path integration may surface issues with
        long-running CI dispatch (LLM costs in CI, especially). Project
        configuration for CI-friendly dispatch (no LLM, longer ATP
        budgets) addresses most of this.
\end{itemize}

\subsection{Exit criteria}

The system is usable for real Lean development. A medium-sized mathlib
file with several dispatch calls compiles in CI within reasonable
time. The dashboard provides actionable feedback to users. Tier 2
either works at acceptable rates or has been formally demoted with
documentation.

\section{Cross-phase concerns}

\subsection{Soundness audit cadence}

The architecture has explicit soundness requirements (translation
discipline, refinement record completeness, rewrite pass soundness
preservation, theory tag soundness). These should be audited at the
end of each phase by someone other than the implementer of the
component being audited. The audit is informal but documented;
findings are tracked.

\subsection{Performance budgets}

Initial performance budgets, to be revised as data accumulates:

\begin{itemize}
  \item Slow-path dispatch (interactive): 30 seconds total budget,
        soft target 5 seconds for typical goals.
  \item Build-path dispatch: 5 minutes total budget per call,
        soft target 30 seconds for typical goals.
  \item IR serialization: under 100ms for typical goals.
  \item IR rewriting: under 1 second for typical goals; pass-level
        timeouts at 5 seconds.
  \item Lifting: under 1 second for typical goals.
\end{itemize}

Budgets that are routinely exceeded should be diagnosed; budgets that
are routinely met by large margins should be tightened.

\subsection{Test coverage strategy}

\begin{itemize}
  \item \textbf{Unit tests} for each component: schema validation,
        per-pass rewrite correctness, per-strategy verifier behavior.
  \item \textbf{Integration tests} for each (home system, backend, tier)
        combination: end-to-end dispatch on canonical goals.
  \item \textbf{Property tests} for round-trip invariants: serialize $\to$
        deserialize, rewrite $\to$ invert, dispatch $\to$ lift.
  \item \textbf{Regression tests} on real-world goals: a curated set
        of mathlib-style goals that the system has historically handled
        correctly. Any regression here is a release blocker.
\end{itemize}

\subsection{Versioning policy}

\begin{itemize}
  \item Spec versions follow semantic versioning. Breaking changes to
        IR schema or certificate format require major version bumps.
  \item The IR \texttt{ir\_version} field and certificate
        \texttt{cert\_version} field track the spec version they
        conform to. Components reject documents with incompatible
        versions.
  \item The patterns registry has its own version, separate from the
        spec. Adding patterns is non-breaking; removing or renaming
        patterns requires a registry version bump.
\end{itemize}

\subsection{Risk register summary}

\begin{tabularx}{\textwidth}{@{}llXl@{}}
\toprule
\textbf{Phase} & \textbf{Risk} & \textbf{Description} & \textbf{Severity} \\
\midrule
1 & Spec gaps & Translation hints insufficient for real Lean code & \highrisk \\
1 & Format quirks & CVC5 certificate format version drift & \medrisk \\
1 & Lean API & Metaprogramming complexity for serializer & \medrisk \\
2 & Replayer scope & Alethe-to-Lean coverage substantial work & \highrisk \\
2 & Inversion bugs & Quotient elimination inversion correctness & \highrisk \\
2 & Embedding lemmas & New mathlib lemmas needed for theory tags & \medrisk \\
3 & LLM ops & Endpoint latency, cost, availability surface late & \medrisk \\
3 & Tuning & Concurrent dispatch grace window calibration & \medrisk \\
4 & Schema fit & Rocq surfaces IR gaps & \highrisk \\
4 & Plugin cost & Rocq plugin engineering effort & \medrisk \\
5 & Tier 2 rate & Reconstruction success rate triggers demotion & \medrisk \\
\bottomrule
\end{tabularx}

\section{What this roadmap does not address}

This roadmap is a technical sequencing document. It does not address:

\begin{itemize}
  \item \textbf{Staffing.} Phase durations assume a small team
        (3--5 engineers) with relevant expertise. Different team sizes
        and skill mixes will require different sequencing.
  \item \textbf{Funding and sustainability.} The architecture is a
        long-term commitment; funding to maintain it across home-system
        and backend version evolution is required.
  \item \textbf{Community engagement.} The Lean and Rocq communities
        will need to be brought along; their feedback should drive
        registry maintenance and prioritization. Community engagement
        process is not specified.
  \item \textbf{Productionization.} Hosted vs.\ self-hosted deployment,
        SLA, security review --- all out of scope for this roadmap.
\end{itemize}

\section{Decision points and revision triggers}

The roadmap has three explicit decision points where progress should
be reassessed:

\paragraph{Phase 1 exit.} Has the IR/spec held up under real implementation?
If yes, proceed to Phase 2. If no, revise spec to v1.1 and re-plan
Phase 2 against the revised spec.

\paragraph{Phase 4 exit.} Did the Rocq probe validate the architecture?
If yes, the architecture is genuinely cross-system and v2 work on
production Rocq integration is sound. If no, the architecture has
Lean-shaped assumptions that need to be re-evaluated --- this could
require significant re-work.

\paragraph{Phase 5 exit / pre-release.} Does the system meet usability
goals? Tier 2 demotion clause, dashboard quality, and CI integration
are the primary gates. A failed pre-release gate triggers a v1.5
focused on the failing area before declaring v1.0 complete.

Beyond these explicit decision points, ongoing assessment of risk
items in the risk register may trigger re-planning at any point.
The roadmap is a plan, not a contract.


\newpage
\section*{Appendix (v1.1, 2026-09-05): the R-series and the ``Phase 6'' double meaning}
\addcontentsline{toc}{section}{Appendix (v1.1): the R-series}

This roadmap is the v1.0 plan and is kept unchanged above. Work after
the return from hiatus is sequenced as an \emph{R-series} of gated
phases, each ending at a checkpoint adjudicated by an adversarial
review until convergence. The R-series roadmap is
\texttt{spec/roadmap-v1.1.md}; the spec v1.1 delta it accompanies is
\texttt{delta.md~\S7}, and the per-phase decision records are
\texttt{delta.md~\S5}. The mapping from the phases above to the
R-series:

\begin{tabularx}{\textwidth}{@{}l>{\hsize=1.25\hsize}X>{\hsize=0.75\hsize}X@{}}
\toprule
\textbf{R-phase} & \textbf{Goal} & \textbf{Absorbs} \\
\midrule
R0 & re-green CI, local parity, repository hygiene, script-derived status & --- \\
R1 & close the Alethe walker's production path (mint gate $=$ walker rule set; UF/UFLIA routed to the walker; live-strict corpus suites) & Phase 2.6 replayer; ``Phase 6'' (walker scale) \\
R2 & make the certificate load-bearing: rewriter on the live path, real \texttt{rewrite\_trace\_hash}, honest envelope, identity-trace guard & Phase 2.4/2.5 \\
R3 & specialization and lifting: $\mathbb{N}\to\mathbb{Z}$, polymorphic $\alpha$, definitional unfolding inverted in the lifted term & Phase 2.1--2.3, 2.7 \\
R4 & external demo: \texttt{by proof\_broker} on a downstream Lake project's obligations & Phase 3 exit criterion, applied externally \\
R5 & spec v1.1 delta, roadmap v1.1, documentation consolidation & the Phase 1 spec-revision checkpoint; the Phase 5 exit gates (Section~\ref{sec:phase5}), disposed below \\
\bottomrule
\end{tabularx}

\paragraph{``Phase 6''.} This document defines Phases 0--5. Two later
bodies of work were both called ``Phase 6'' in the repository: the
cross-platform distribution work moved out of Phase 5 (the CI matrix,
the install layout, the macOS signing scaffold --- ``Phase 6-D''), and
the walker scale profile (``Phase 6-W'', absorbed by R1). A single
``Phase 7'' remark in a retrospective was a hypothetical, never a
phase. The v1.0 numbering is frozen as history; new work is
R-numbered. Details, gates, and the decide-list that follows R5:
\texttt{spec/roadmap-v1.1.md}.

\paragraph{Phase 5 exit, as it turned out.} Tier~2's lemma-list form
is formally demoted with documentation (\texttt{delta.md~\S7.4(b)});
the dashboard and the build-path CI integration were not built and
are decide-list items; there is no single ``v1.0 complete''
declaration --- the R-series gates are the release criteria, phase by
phase.

\end{document}
